% =====================================================================
%  DSM1 — Sistema Multiagente de Controladoria Terceirizada
%  INE5628 — Sistemas Multiagentes (UFSC/CTC/INE)
%
%  Compilar:  latexmk -pdf main.tex     (ou: make pdf, na raiz)
%  As tabelas em generated/ são geradas por tools/trace.py a partir de
%  models/*.yaml — não edite à mão; edite os YAML e rode `make trace`.
% =====================================================================
\documentclass[10pt,a4paper]{article}

% ------------------------------------------------------------ autoria
\newcommand{\autorA}{Danrley Moreira Ribeiro}
\newcommand{\matA}{23102664}
\newcommand{\autorB}{Jiliard Mai Peifer}
\newcommand{\matB}{23103129}
\newcommand{\repoURL}{https://github.com/danrley-ribeiro/fin-dept-multiagents-project}
\newcommand{\docsURL}{https://fin-dept-multiagents-project.readthedocs.io}
\newcommand{\versao}{1.0}

% ------------------------------------------------------------ pacotes
\usepackage[T1]{fontenc}
\usepackage[utf8]{inputenc}
\usepackage[brazil]{babel}
\usepackage[default,semibold]{sourcesanspro}
\usepackage[a4paper, left=1.75cm, right=1.75cm, top=2.15cm, bottom=1.9cm, headsep=0.55cm, footskip=0.9cm]{geometry}
\usepackage[table]{xcolor}
\usepackage{amsmath, amssymb}
\usepackage{microtype}
\usepackage{array, booktabs, xltabular, makecell}
\usepackage{enumitem}
\usepackage{titlesec}
\usepackage{fancyhdr}
\usepackage{ragged2e}
\usepackage{graphicx}
\usepackage{caption}
\usepackage{float}
\usepackage[section]{placeins}
\usepackage{needspace}
\usepackage[most]{tcolorbox}
\usepackage{tikz}
\usetikzlibrary{arrows.meta, positioning, shapes.geometric, shapes.misc, calc, fit, backgrounds}
\usepackage{newunicodechar}
\newunicodechar{⟨}{\ensuremath{\langle}}
\newunicodechar{⟩}{\ensuremath{\rangle}}
\newunicodechar{≠}{\ensuremath{\neq}}
\newunicodechar{→}{\ensuremath{\rightarrow}}
\usepackage{csquotes}
\usepackage[backend=bibtex, style=authoryear, maxcitenames=2, uniquelist=false, dashed=false, giveninits=true]{biblatex}
\addbibresource{referencias.bib}
\DeclareFieldFormat{urldate}{\mkbibparens{acesso em #1}}
\appto\biburlsetup{\Urlmuskip=0mu\relax}
\renewcommand*{\nameyeardelim}{\addcomma\space}
\usepackage{titletoc}
\usepackage[hidelinks, pdfusetitle]{hyperref}
\hypersetup{pdftitle={DSM1 — Sistema Multiagente de Controladoria Terceirizada},
            pdfauthor={\autorA, \autorB}, pdfsubject={INE5628 — Entrega 1}}

% ------------------------------------------------------------ cores
\definecolor{navy}{HTML}{111A2E}
\definecolor{accent}{HTML}{1565A8}
\definecolor{accentlight}{HTML}{E7F0FA}
\definecolor{muted}{HTML}{6B7280}
\definecolor{rowalt}{HTML}{F5F7FA}
\definecolor{plat}{HTML}{5B4B8A}
\definecolor{platlight}{HTML}{EFEBF7}
\definecolor{warn}{HTML}{B45309}
\definecolor{bad}{HTML}{B91C1C}
\definecolor{good}{HTML}{15803D}

% ------------------------------------------------------------ sumário
\titlecontents{section}[2.2em]{\addvspace{25pt}\large\bfseries}{\contentslabel{2.2em}}{}{\titlerule*[0.7pc]{.}\contentspage}[\addvspace{4pt}]
\titlecontents{subsection}[4.9em]{\addvspace{10pt}\normalsize}{\contentslabel{2.7em}}{}{\titlerule*[0.7pc]{.}\contentspage}
\setcounter{tocdepth}{2}

% ------------------------------------------------------------ tipografia e layout
\setlength{\parindent}{0pt}
\setlength{\parskip}{0.45em}
\linespread{1.04}
\renewcommand{\arraystretch}{1.18}
\setlist{nosep, leftmargin=1.1em, itemsep=0.15em, topsep=0.2em}
\captionsetup{font={small,sf}, labelfont=bf, skip=4pt}

\titleformat{\section}{\large\bfseries}{\thesection.}{0.5em}{\MakeUppercase}
\titlespacing*{\section}{0pt}{1.3em}{0.6em}
\titleformat{\subsection}{\normalsize\bfseries}{\thesubsection}{0.5em}{}
\titlespacing*{\subsection}{0pt}{0.7em}{0.25em}

\pagestyle{fancy}
\fancyhf{}
\renewcommand{\headrulewidth}{0pt}
\fancyhead[L]{\scriptsize INE5628 / DSM1 — Sistema Multiagente de Controladoria Terceirizada}
\fancyhead[R]{\scriptsize Entrega 1: Cenário, Requisitos e Modelagem}
\fancyfoot[L]{\scriptsize UFSC / CTC / INE — Prof. Rafael Luiz Cancian}
\fancyfoot[R]{\scriptsize Página \thepage}
\fancypagestyle{first}{\fancyhf{}\renewcommand{\headrulewidth}{0pt}
  \fancyfoot[L]{\scriptsize UFSC / CTC / INE — Sistemas Multiagentes}
  \fancyfoot[R]{\scriptsize Página \thepage}}

% ------------------------------------------------------------ caixas
\newenvironment{callout}{\par\addvspace{0.6em}\begingroup}{\endgroup\par\addvspace{0.6em}}
\newcommand{\callouttitle}[1]{{\bfseries #1}\par\vspace{2pt}}

% ------------------------------------------------------------ tabelas
\newcolumntype{L}[1]{>{\RaggedRight\arraybackslash}p{#1}}
\newcolumntype{Y}{>{\RaggedRight\arraybackslash}X}
\newcommand{\hd}[1]{\textbf{#1}}
\newcommand{\headrow}{\toprule}
\newcommand{\rowsep}{\specialrule{0.25pt}{1.5pt}{1.5pt}}
\setlength{\heavyrulewidth}{0.5pt}\setlength{\lightrulewidth}{0.3pt}
\newenvironment{dtable}[1]{\footnotesize\renewcommand{\arraystretch}{1.3}\xltabular{\linewidth}{#1}}{\bottomrule\endxltabular\par\addvspace{0.4em}}

% ------------------------------------------------------------ rastreabilidade
\newcommand{\idtarget}[1]{\leavevmode\textbf{#1}\hypertarget{#1}{}}
\newcommand{\idref}[1]{\hyperlink{#1}{#1}}
\newcommand{\llmflag}{\textbf{flag}}
\newcommand{\llmno}{não}
\newcommand{\decisao}[1]{\textbf{\texttt{\detokenize{#1}}}}
\newcommand{\code}[1]{\texttt{\small #1}}

% ARQUIVO GERADO por tools/trace.py a partir de models/*.yaml — não editar à mão
\newcommand{\nReq}{18}
\newcommand{\nObj}{136}
\newcommand{\nTest}{25}
\newcommand{\nMet}{12}
\newcommand{\nProt}{6}
\newcommand{\nNorm}{8}
\newcommand{\nPend}{38}


% =====================================================================
\begin{document}
% ------------------------------------------------------------ capa (página exclusiva)
\begin{titlepage}
\centering
{\large\bfseries UNIVERSIDADE FEDERAL DE SANTA CATARINA\par}
{\normalsize Centro Tecnológico \textbullet{} Departamento de Informática e Estatística\par}
{\normalsize INE5628 — Sistemas Multiagentes\par}
\vfill
{\fontsize{30}{36}\selectfont\bfseries DSM1\par}
\vspace{0.9cm}
{\fontsize{22}{28}\selectfont\bfseries Sistema Multiagente de\\Controladoria Terceirizada\par}
\vspace{0.7cm}
{\Large Entrega 1: Cenário-Problema, Requisitos,\\Modelagem AOSE e Rastreabilidade\par}
\vfill
{\large\bfseries Autores\par}\vspace{0.3cm}
{\large \autorA{} (\matA)\par \autorB{} (\matB)\par}
\vspace{1cm}
{\large\bfseries Professor\par}\vspace{0.3cm}
{\large Prof. Dr. Rafael Luiz Cancian\par}
\vfill
{\normalsize Peso: 20\% da MPDSM \textbullet{} Prazo: 10/10/2026 \textbullet{} Versão \versao\par}
\vspace{0.3cm}
{\normalsize Documentação: \href{\docsURL}{\texttt{\docsURL}}\par}
\vspace{0.8cm}
{\large Florianópolis, 2026\par}
\end{titlepage}

% ------------------------------------------------------------ sumário (página exclusiva)
\thispagestyle{empty}
{\centering\fontsize{20}{24}\selectfont\bfseries Sumário\par}
\vspace{1.2cm}
\makeatletter\@starttoc{toc}\makeatother
\clearpage
\setcounter{page}{1}

\begin{callout}
\callouttitle{Resumo executivo do projeto}
Este documento especifica a Entrega 1 (DSM1) de um \textbf{Sistema Multiagente (SMA) para terceirização da
Controladoria}: um orquestrador e cinco agentes especializados (Custos e Rentabilidade, Orçamento, Resultados e
Indicadores, Controles Internos e Relatórios Gerenciais), apoiados por serviços de plataforma de comunicação
(A2A/MCP), auditoria, configuração de modelos, memória e política de segurança. A modelagem segue AOSE \parencite{cancian2026},
combinando Tropos \parencite{bresciani2004tropos}, Gaia \parencite{wooldridge2000gaia} e Prometheus
\parencite{padgham2004prometheus}, e estabelece a justificativa do SMA, a fronteira e o ambiente,
\nReq{} requisitos com evidência observável, papéis $\langle$Resp, Perm, Obl, Prot$\rangle$, \nProt{} protocolos
como máquinas de estados, \nNorm{} normas e \nMet{} métricas. A IA generativa é \textbf{opcional e desligada por
padrão}, restrita à narrativa em três agentes. Todos os \nObj{} objetos possuem ID estável e formam uma cadeia de
rastreabilidade validada automaticamente, publicada em \href{\docsURL}{documentação navegável}.
\end{callout}


% =====================================================================
\section{Cenário-problema, motivação e objetivos}

\subsection{Contexto e problema operacional}
Uma empresa de médio porte terceiriza (BPO) sua Controladoria. Todo mês é preciso fechar o resultado
(balancete $\rightarrow$ rateios $\rightarrow$ DRE $\rightarrow$ indicadores $\rightarrow$ relatório) e, a cada
trimestre, avaliar o orçamento. Hoje o ciclo leva dias de consolidação manual (\code{PROB-01}), não mostra margem
por produto ou unidade (\code{PROB-02}), detecta desvios orçamentários tarde (\code{PROB-03}), deixa passar riscos de
alçada, segregação de funções e duplicidade (\code{PROB-04}) e gera relatórios inconsistentes com a DRE
(\code{PROB-05}). Automatizar com IA traz problemas novos: decisões sem trilha nem controle de custo
(\code{PROB-06}), conhecimento que se perde entre ciclos (\code{PROB-07}) e ações com efeito externo sem
controle humano (\code{PROB-08}).

\textbf{Stakeholders} (recorte Tropos; \cite{bresciani2004tropos}): CFO/diretoria contratante, controller humano (aprovador), contabilidade
do cliente, auditoria interna/externa, TI/Segurança e DPO (LGPD), equipe de desenvolvimento e avaliador.
\textbf{Cenários âncora (MVP):} \emph{“feche o resultado do mês”} e \emph{“avalie o orçamento do trimestre”},
sobre uma empresa sintética com gabarito, com fontes acessadas por \emph{MCP servers} simulados e offline.

\subsection{Justificativa para o paradigma multiagente (hipótese MAS)}
Aplicando o checklist de adequação de \textcite[Tab.~5.1]{cancian2026}:

\begin{dtable}{L{2.25cm} Y L{1.05cm}}
\headrow \hd{Critério} & \hd{Evidência no domínio} & \hd{MAS?}\\ \midrule\endhead
Autonomia & Controles Internos decide sozinho \textbf{vetar} a publicação; Orçamento decide o que é desvio relevante pelos próprios limiares. & \textbf{sim}\\ \rowsep
Distribuição & Dados em fontes heterogêneas (ERP, planilhas, histórico, mercado, normas); nenhum agente tem o estado global. & \textbf{sim}\\ \rowsep
Interdependência & Resultados depende de Custos e Orçamento; o veto de Controles Internos altera as opções do Orquestrador e de Relatórios. & \textbf{sim}\\ \rowsep
Heterogeneidade & Cálculo determinístico (Custos), regras BDI (Controles), síntese com LLM opcional (Relatórios). & \textbf{sim}\\ \rowsep
Recursos & Período contábil exclusivo durante o fechamento; orçamento de tokens consumível. & \textbf{sim}\\ \rowsep
Evidência & Cenários com gabarito e protocolos com transições inválidas tornam o comportamento testável. & \textbf{sim}\\
\end{dtable}

\textbf{Condição explícita de rejeição.} Se a demanda se reduzir a gerar um único relatório por uma sequência fixa
de etapas, sem veto, sem dependência entre unidades de decisão e sem recurso disputado, um \emph{workflow}
determinístico basta. Esse workflow sequencial é a \textbf{baseline} das entregas seguintes. \textbf{Hipóteses
verificáveis:} (H1) o veto de Controles Internos barra 100\% das inconformidades críticas plantadas antes da
publicação; (H2) a decomposição mantém 100\% de acurácia contra o gabarito, com custo de comunicação mensurável;
(H3) com a flag LLM desligada, todos os relatórios são gerados sem rede e a custo zero.

\subsection{Objetivos do sistema}
\textbf{Objetivo geral (\code{G-00}):} executar as rotinas de controladoria de forma terceirizada e coordenada,
entregando à liderança o resultado do período e análises confiáveis, respeitando alçadas, sigilo e aprovação
humana, com evidência reconstruível de cada decisão. \textbf{Objetivos específicos:} (G-01) atender solicitações
decompondo-as em tarefas coordenadas; (G-02) apurar custos, margens e rentabilidade; (G-03) controlar o orçamento
e projetar forecast; (G-04) consolidar DRE, EBITDA, fluxo de caixa e KPIs; (G-05) garantir conformidade antes da
publicação; (G-06) entregar relatórios executivos acionáveis; (G-07) manter auditabilidade total e custo de IA sob
controle; (G-08) reutilizar conhecimento aprovado; (G-09) restringir efeitos externos à autorização humana.

% =====================================================================
\section{Fronteira do sistema e caracterização do ambiente}
Pertencem ao \textbf{software modelado} os seis agentes e os sete serviços de plataforma (Figura~\ref{fig:org}).
Permanecem no \textbf{ambiente}, ainda que representados por adapters, os usuários humanos (solicitante, aprovador,
auditor), as fontes externas e o relógio/calendário de fechamento.

\begin{figure}[H]
  \centering
  \resizebox{\linewidth}{!}{% Figura editável (TikZ): organização em camadas do SMA de Controladoria (DSM1).
% IDs conforme models/*.yaml. Incluída por main.tex via \input.
\begin{tikzpicture}[
  font=\sffamily\scriptsize, >={Stealth[length=2mm]},
  band/.style={draw=accent!35, fill=accentlight, rounded corners=4pt},
  bandlabel/.style={font=\sffamily\bfseries\tiny, text=accent, anchor=north west},
  ag/.style={draw=navy, fill=white, rounded corners=2pt, minimum width=2.05cm, minimum height=0.95cm, align=center, line width=0.6pt},
  svc/.style={draw=plat, fill=platlight, rounded corners=2pt, minimum width=3.3cm, minimum height=0.62cm, align=center},
  comm/.style={draw=accent, fill=white, rounded corners=2pt, minimum height=0.6cm, align=center},
  src/.style={cylinder, shape border rotate=90, aspect=0.2, draw=muted, fill=gray!8, minimum width=1.75cm, minimum height=0.75cm, align=center},
  hum/.style={draw=navy, fill=navy!7, rounded corners=7pt, minimum height=0.62cm, align=center},
  llm/.style={font=\sffamily\bfseries\tiny, text=warn},
  lbl/.style={font=\sffamily\tiny, fill=white, inner sep=1pt, text=navy},
]
% ---------- humanos (fora da fronteira)
\node[hum] (sol) at (6.0,7.55) {\textbf{Solicitante} (CFO / Diretoria) · ROLE-SOL};
\node[hum] (apr) at (14.55,7.55) {\textbf{Aprovador} (controller) · ROLE-APR};

% ---------- fronteira do sistema
\draw[navy, dashed, rounded corners=6pt, line width=0.7pt] (-0.2,1.95) rectangle (16.45,7.0);
\node[font=\sffamily\bfseries\tiny, text=navy, anchor=south west] at (-0.15,7.0) {FRONTEIRA DO SMA};

% ---------- camada de orquestração
\draw[band] (0.0,5.75) rectangle (12.1,6.85);
\node[bandlabel] at (0.05,6.85) {ORQUESTRAÇÃO};
\node[ag, minimum width=3.6cm, fill=navy, text=white] (orq) at (6.0,6.25) {\textbf{AG-ORQ · Orquestrador}\\[-1pt]{\tiny decompõe, delega, agrega}};

% ---------- agentes especializados
\draw[band] (0.0,3.75) rectangle (12.1,5.45);
\node[bandlabel] at (0.05,5.45) {AGENTES ESPECIALIZADOS};
\node[ag] (cus) at (1.25,4.5) {\textbf{AG-CUS}\\Custos e\\Rentabilidade};
\node[ag] (res) at (3.65,4.5) {\textbf{AG-RES}\\Resultados e\\Indicadores};
\node[ag] (orc) at (6.05,4.5) {\textbf{AG-ORC}\\Orçamento e\\Controle};
\node[ag] (cin) at (8.45,4.5) {\textbf{AG-CIN}\\Controles\\Internos};
\node[ag] (rel) at (10.85,4.5) {\textbf{AG-REL}\\Relatórios\\Gerenciais};
\foreach \n in {res,orc,rel} \node[llm, anchor=north] at (\n.south) {LLM opcional (flag)};

% ---------- comunicação
\draw[band] (0.0,2.1) rectangle (12.1,3.25);
\node[bandlabel] at (0.05,3.25) {COMUNICAÇÃO};
\node[comm, minimum width=7.1cm] (a2a) at (4.25,2.6) {\textbf{SVC-A2A} · agente $\leftrightarrow$ agente (FIPA-ACL, conversation\_id)};
\node[comm, minimum width=3.6cm] (mcp) at (10.0,2.6) {\textbf{SVC-MCP} · agente $\rightarrow$ fonte (tools)};

% ---------- serviços de plataforma (transversais)
\draw[band, draw=plat!50, fill=platlight!40] (12.45,2.1) rectangle (16.3,6.85);
\node[bandlabel, text=plat] at (12.5,6.85) {SERVIÇOS DE PLATAFORMA};
\node[svc] (pol) at (14.37,6.15) {\textbf{SVC-POL} SafetyPolicy / Approval};
\node[svc] (aud) at (14.37,5.35) {\textbf{SVC-AUD} Auditoria (JSONL)};
\node[svc] (mod) at (14.37,4.55) {\textbf{SVC-MOD} Registro de modelos};
\node[svc] (mem) at (14.37,3.75) {\textbf{SVC-MEM} Memória (emb.\ + grafo)};
\node[svc] (lck) at (14.37,2.95) {\textbf{SVC-LCK} Reserva de período};

% ---------- ambiente (fontes externas)
\node[font=\sffamily\bfseries\tiny, text=muted, anchor=west] at (-0.15,1.45) {AMBIENTE};
\node[src] (erp) at (1.4,0.85) {ERP};
\node[src] (pln) at (3.8,0.85) {Orçamento};
\node[src] (bd)  at (6.2,0.85) {Histórico};
\node[src] (api) at (8.6,0.85) {Mercado};
\node[src] (nrm) at (11.0,0.85) {Normas};

% ---------- interações
\draw[->, navy, thick] (sol) -- node[lbl, right] {PROT-SOL} (orq);
\draw[->, navy, thick] (orq.south) -- node[lbl] {PROT-DEL · PROT-CMP (veto)} (6.05,5.45);
\draw[->, accent, thick] (res) -- node[lbl, above] {QRY} (cus);
\draw[->, accent, thick] (res.north east) to[bend left=18] node[lbl, above] {QRY} (orc.north west);
\foreach \n in {cus,res,orc,cin,rel} \draw[accent!60, line width=0.6pt] (\n.south |- 0,3.75) -- (\n.south |- 0,3.25);
\draw[->, muted, thick] (mcp.south) -- (10.0,1.55) -| (erp.north);
\foreach \n in {pln,bd,api,nrm} \draw[->, muted] (\n.north |- 0,1.55) -- (\n.north);
\draw[->, plat, thick] (orq.east) -- node[lbl, above] {PROT-LCK · PROT-APR} (pol.west);
\draw[<->, plat, thick] (apr) -- node[lbl, right] {aprovação / waiver} (pol);
\draw[->, plat, dashed] (12.1,4.55) -- node[lbl, below, font=\sffamily\tiny] {intercepta} (mod.west);
\end{tikzpicture}
}
  \caption{Organização em camadas: fronteira do SMA, agentes, serviços de plataforma transversais, protocolos e ambiente.}
  \label{fig:org}
\end{figure}

O ambiente é:
\begin{itemize}
  \item \textbf{Parcialmente observável:} cada agente vê apenas o que seus conectores MCP e mensagens A2A entregam
        ($see_i: S \to P_i$): Custos não lê a base normativa e Controles Internos não lê as regras de rateio.
  \item \textbf{Dinâmico e estocástico:} lançamentos e estornos chegam durante o fechamento ($next: S\times A\to S$
        muda também por eventos externos); APIs falham ou atrasam; a resposta humana tem duração incerta.
  \item \textbf{Discreto e sequencial por ciclo:} contas, competências e lançamentos têm identificadores bem
        definidos; ajustes aprovados e memória de um ciclo afetam o próximo.
  \item \textbf{Com recursos compartilhados:} o período em fechamento é exclusivo e o orçamento de tokens é consumível.
\end{itemize}


Recursos externos classificados por observabilidade, acionabilidade, exclusividade e autoridade
(escrita em fonte externa é sempre ação restrita):

\begin{dtable}{L{3.4cm} Y Y L{2.6cm} L{2.6cm}}
\headrow \hd{Recurso} & \hd{Observabilidade} & \hd{Acionabilidade} & \hd{Exclusividade} & \hd{Autoridade}\\ \midrule\endhead
% ARQUIVO GERADO por tools/trace.py a partir de models/*.yaml — não editar à mão
\textbf{ERP — razão contábil, balancete e lançamentos} & Leitura via MCP (balancete, lançamentos, centros de custo, plano de contas) & Escrita de ajuste contábil — ação restrita (REQUIRE\_APPROVAL) & Período contábil exclusivo durante o fechamento & Contabilidade do cliente (STK-03) \\
\textbf{Planilhas de orçamento e premissas} & Leitura via MCP (orçado por conta/centro de custo, premissas) & Alteração de orçamento — ação restrita (REQUIRE\_APPROVAL) & Não exclusiva para leitura & Diretoria / FP\&A do cliente (STK-01) \\
\textbf{Base analítica histórica (data warehouse)} & Leitura via MCP (séries históricas, fechamentos anteriores) & Nenhuma (somente leitura) & Não exclusiva & TI do cliente (STK-05) \\
\textbf{APIs financeiras de mercado (câmbio, IPCA, Selic)} & Leitura via MCP; pode falhar ou atrasar & Nenhuma & Não exclusiva (cota de chamadas) & Provedor externo \\
\textbf{Base normativa interna (políticas de alçada, SoD, manual contábil)} & Leitura via MCP/RAG com proveniência (documento, trecho, versão) & Nenhuma & Não exclusiva & Compliance / Auditoria (STK-04) \\
\textbf{Período contábil em fechamento (recurso compartilhado)} & Estado de reserva consultável & Reservar / liberar via protocolo de reserva & Exclusivo por (entidade, competência), com TTL & SVC-LCK \\
\textbf{Orçamento de tokens e custo de IA (recurso consumível)} & Saldo por agente/tarefa consultável no registro de modelos & Consumido apenas por chamadas LLM habilitadas por flag & Cota por agente e por tarefa & Controller (STK-02) via SVC-MOD \\
\textbf{Interface do usuário (solicitação, aprovação e observabilidade)} & Recebe solicitações e decisões de aprovação humanas & Exibe resultado final, pedidos de aprovação e dashboard & Não exclusiva & Usuários autenticados (STK-01, STK-02, STK-04) \\

\end{dtable}

% =====================================================================
\section{Especificação de requisitos e critérios de aceite}
Os requisitos seguem o template AOSE \parencite[Cap.~4]{cancian2026}: cada um liga-se a um papel responsável e a uma
\textbf{evidência observável} capaz de distinguir implementação correta de execução apenas aparente. Os IDs
reaparecem nos testes derivados, nas métricas e na matriz da Seção~\ref{sec:trace}.

\begin{dtable}{L{1.5cm} L{1.95cm} Y L{1.45cm} L{4.4cm}}
\headrow \hd{ID} & \hd{Tipo} & \hd{Requisito} & \hd{Papéis} & \hd{Evidência observável de aceite}\\ \midrule\endhead
% ARQUIVO GERADO por tools/trace.py a partir de models/*.yaml — não editar à mão
\idtarget{REQ-F01} & Funcional & \textbf{Decomposição e agregação de solicitações.} O Orquestrador deve mapear cada solicitação do catálogo ("fechar o resultado do mês", "avaliar o orçamento do trimestre") para um DAG de subtarefas, delegá-las aos papéis competentes e agregar uma resposta final; pedido fora do catálogo recebe not-understood explícito. & ORQ, SOL & Trace com o DAG gerado e, por subtarefa, request $\rightarrow$ agree $\rightarrow$ inform/failure sob o mesmo conversation\_id; resposta final com status explícito. \\
\idtarget{REQ-F02} & Funcional & \textbf{Apuração de custos, margem e rentabilidade.} Apurar custos fixos e variáveis, aplicar rateio por critério configurado e calcular margem de contribuição e rentabilidade por produto, serviço e UN a partir do balancete obtido via MCP. & CUS & Valores idênticos ao gabarito do dataset sintético (tolerância R\$ 0,01) e cada valor acompanhado de proveniência (fonte MCP, consulta, versão do rateio). \\
\idtarget{REQ-F03} & Funcional & \textbf{Orçado x realizado e forecast.} Comparar orçado x realizado por conta e centro de custo, sinalizar desvios acima de limiar configurável (percentual e absoluto) e projetar o forecast do trimestre com método declarado. & ORC & Desvios sinalizados exatamente nos itens plantados acima do limiar do cenário de teste; forecast reproduzível com o mesmo seed. \\
\idtarget{REQ-F04} & Funcional & \textbf{Resultado consolidado via A2A.} Montar DRE gerencial, EBITDA, fluxo de caixa (método indireto) e KPIs, obtendo custos e orçamento dos agentes especialistas por query-ref A2A, sem reimplementar suas regras. & RES, CUS, ORC & Trace mostra query-ref RES$\rightarrow$CUS e RES$\rightarrow$ORC antes do inform final; DRE igual ao gabarito; inform não solicitado é rejeitado. \\
\idtarget{REQ-F05} & Funcional & \textbf{Controles internos com veto.} Verificar alçada, segregação de funções, duplicidade, conciliação pendente e competência sobre o pacote do fechamento, emitindo parecer com regra violada e severidade; não conformidade crítica veta a publicação. & CIN, ORQ & Todas as inconformidades críticas plantadas detectadas; publicação em estado VETADO é bloqueada até correção ou waiver humano. \\
\idtarget{REQ-F06} & Funcional & \textbf{Relatório executivo por template com LLM opcional.} Gerar relatório executivo por template determinístico a partir dos resultados consolidados; com a flag LLM ativa, acrescentar narrativa e insights preservando os números do núcleo. & REL & Com flag desligada o relatório é gerado sem rede; com flag ligada todo número do texto coincide com o núcleo ou o template é usado. \\
\idtarget{REQ-F07} & Funcional & \textbf{Memória de longo prazo com proveniência.} Toda saída aprovada pelo controller vira memória (embedding + nós e arestas de conta, indicador, decisão e período); os agentes consultam a memória antes de nova análise e citam a proveniência usada. & CUS, ORC, RES, REL & Saída não aprovada não aparece na memória; análise seguinte cita o ID da memória e a fonte original. \\
\idtarget{REQ-S01} & Segurança & \textbf{Ações com efeito externo só com aprovação.} Lançar ajuste no ERP, alterar orçamento ou publicar relatório nunca é executado autonomamente: a SafetyPolicy classifica em ALLOW, REQUIRE\_APPROVAL ou DENY e retém a ação no Approval Gate. & ORQ, REL, ORC, APR & ActionRequest em PENDENTE até approve humano; ação DENY nunca executa efeito colateral, mesmo com aprovação. \\
\idtarget{REQ-S02} & Segurança & \textbf{Exclusão mútua do período em fechamento.} Dois fechamentos concorrentes da mesma entidade e competência são impedidos por reserva exclusiva com TTL; o conflito retorna falha explicada. & ORQ & Segunda reserva sobreposta recebe reserve-failure com o dono atual; reserva expirada é liberada. \\
\idtarget{REQ-S03} & Segurança & \textbf{LLM não calcula nem autoriza.} Todo número presente em texto gerado por LLM é validado contra o núcleo determinístico; divergência descarta o texto e aciona o template. Proposta de ação vinda de LLM nunca é tratada como autorização. & REL, RES, ORC & Mock de LLM que altera um valor provoca fallback registrado em log; mock que sugere ação não gera execução. \\
\idtarget{REQ-O01} & Observabilidade & \textbf{Log de decisão.} Cada decisão registra critério/regra, referências dos dados de entrada, agente, conversation\_id e timestamp em JSONL append-only, via interceptor central (não por agente). & ORQ, CUS, ORC, RES, CIN, REL & Script reconstrói a sequência causal de uma solicitação a partir do conversation\_id sem inspecionar o código. \\
\idtarget{REQ-O02} & Observabilidade & \textbf{Log de prompt.} Quando a flag LLM estiver ativa, registrar o prompt exato (após redaction), versão do template, modelo e fontes anexadas. & REL, RES, ORC & Com flag ligada há um registro de prompt por chamada; com flag desligada nenhum registro de prompt é criado. \\
\idtarget{REQ-O03} & Observabilidade & \textbf{Log de custo e tokens.} Registrar tokens de entrada e saída, modelo e custo estimado por chamada, agregados por tarefa, solicitação e agente. & REL, RES, ORC & Soma dos custos por chamada igual ao total da tarefa; custo zero com flag desligada. \\
\idtarget{REQ-O04} & Observabilidade & \textbf{Interface de observabilidade segura.} Dashboard web, mobile e desktop autenticado com filtros por agente, período, tarefa e custo sobre os logs de auditoria. Especificado na DSM1; implementação em entrega futura. & APR, SOL & Cada filtro retorna exatamente os registros do recorte em base de teste; acesso sem autenticação é negado. \\
\idtarget{REQ-C01} & Configuração & \textbf{Modelo por agente trocável sem redeploy.} Registro central define por agente a flag LLM (padrão desligada), modelo, limite de tokens e orçamento; a mudança vale na próxima tarefa sem redeploy; LLM só é habilitável em REL, RES e ORC; cota excedida bloqueia a chamada e aciona o template. & REL, RES, ORC, APR & Troca de modelo refletida na tarefa seguinte; configuração com LLM em ORQ, CUS ou CIN é rejeitada; cota estourada gera fallback registrado. \\
\idtarget{REQ-P01} & Privacidade & \textbf{Redaction fail-closed (LGPD e sigilo).} Logs, traces e memória passam por redaction de segredos, credenciais e dados pessoais, com política de retenção; falha de sanitização impede a gravação. & ORQ, CIN, REL & Registro com chave de API ou CPF sintético é sanitizado ou recusado; nenhum campo proibido aparece nos arquivos JSONL. \\
\idtarget{REQ-NF01} & Não funcional & \textbf{Execução offline e reprodutível.} O núcleo executa sem rede e sem API paga, com MCP mocks, dados sintéticos e seed fixo; o LLM é estritamente opcional. & ORQ, REL & Duas execuções com o mesmo seed e flag desligada produzem saídas com o mesmo hash. \\
\idtarget{REQ-NF02} & Não funcional & \textbf{Versionamento de prompts, regras e configuração.} Templates de prompt, regras de conformidade e configurações de modelo têm ID e versão, registrados em cada log de decisão e de prompt. & REL, CIN & Todo registro de log contém as versões vigentes; mudança de versão aparece no trace seguinte. \\

\end{dtable}

% =====================================================================
\section{Agentes, papéis e arquitetura da organização}
O modelo organizacional (Gaia; \cite{wooldridge2000gaia}) separa \textbf{papéis} (o que precisa ser feito, sob quais direitos e deveres),
\textbf{agentes} (entidades que os assumem) e \textbf{serviços de plataforma}. Cada papel é a tupla
$role = \langle Resp, Perm, Obl, Prot\rangle$. Os papéis humanos \code{ROLE-SOL} (solicitante) e \code{ROLE-APR}
(aprovador) ficam fora da fronteira, mas participam de protocolos.

\begin{dtable}{L{2.75cm} Y Y Y L{1.35cm} L{0.75cm}}
\headrow \hd{Agente / cognição} & \hd{Responsabilidades (Resp)} & \hd{Permissões (Perm)} & \hd{Obrigações (Obl)} & \hd{Protocolos} & \hd{LLM}\\ \midrule\endhead
% ARQUIVO GERADO por tools/trace.py a partir de models/*.yaml — não editar à mão
\textbf{Orquestrador (Controladoria Master)}\newline{\scriptsize\color{muted}AG-ORQ · Reativo a protocolos + biblioteca de planos (intenção $\rightarrow$ DAG)} & \textbullet~Receber solicitações e mapear para intenção do catálogo\newline \textbullet~Decompor em DAG de subtarefas e delegar aos papéis competentes\newline \textbullet~Reservar o período e agregar a resposta final & \textbullet~Delegar (request) a qualquer papel especialista\newline \textbullet~Reservar e liberar período no PeriodLockService\newline \textbullet~Propor publicação ao Approval Gate & \textbullet~Responder toda solicitação com resultado ou falha explícita\newline \textbullet~Respeitar reply\_by e veto de conformidade\newline \textbullet~Liberar a reserva ao término ou falha & SOL, DEL, CMP, LCK, APR & \llmno \\
\textbf{Agente de Custos e Rentabilidade}\newline{\scriptsize\color{muted}AG-CUS · Deliberativo determinístico (regras de rateio + cálculo)} & \textbullet~Apurar custos fixos e variáveis e aplicar rateio por critério configurado\newline \textbullet~Calcular margem de contribuição e rentabilidade por produto, serviço e UN & \textbullet~Ler ERP e base histórica via MCP\newline \textbullet~Responder query-ref de outros papéis & \textbullet~Anexar proveniência a cada valor\newline \textbullet~Informar falha se dado de origem estiver ausente & DEL, QRY & \llmno \\
\textbf{Agente de Orçamento e Controle Orçamentário}\newline{\scriptsize\color{muted}AG-ORC · Determinístico (variância + forecast) com narrativa LLM opcional} & \textbullet~Comparar orçado x realizado por conta e centro de custo\newline \textbullet~Sinalizar desvios acima do limiar e projetar forecast & \textbullet~Ler planilhas de orçamento e ERP via MCP\newline \textbullet~Propor alteração de orçamento (restrita)\newline \textbullet~Responder query-ref & \textbullet~Explicitar limiar e método de forecast usados\newline \textbullet~Não alterar orçamento sem aprovação & DEL, QRY, APR & \llmflag \\
\textbf{Agente de Resultados e Indicadores}\newline{\scriptsize\color{muted}AG-RES · Determinístico (DRE, EBITDA, DFC, KPIs) com comentário LLM opcional} & \textbullet~Montar DRE gerencial, EBITDA, fluxo de caixa e KPIs\newline \textbullet~Obter custos e orçamento dos papéis especialistas & \textbullet~Iniciar query-ref para CUS e ORC\newline \textbullet~Ler ERP e APIs de mercado via MCP & \textbullet~Não recalcular rateio próprio — usar o de CUS\newline \textbullet~Sinalizar inconsistência de fechamento & DEL, QRY & \llmflag \\
\textbf{Agente de Controles Internos}\newline{\scriptsize\color{muted}AG-CIN · BDI com regras simbólicas (crenças sobre lançamentos, intenção de vetar)} & \textbullet~Verificar alçada, segregação de funções, duplicidade, conciliação e competência\newline \textbullet~Emitir parecer conforme / não conforme com regra e severidade & \textbullet~Ler ERP e base normativa via MCP\newline \textbullet~Vetar publicação por não conformidade crítica & \textbullet~Citar regra e evidência de cada apontamento\newline \textbullet~Reavaliar após waiver ou correção & DEL, CMP & \llmno \\
\textbf{Agente de Relatórios Gerenciais e Análises}\newline{\scriptsize\color{muted}AG-REL · Template determinístico + LLM opcional (híbrido)} & \textbullet~Sintetizar resultados em relatório executivo por template\newline \textbullet~Acrescentar narrativa e insights quando a flag LLM estiver ativa & \textbullet~Consultar memória aprovada\newline \textbullet~Propor publicação do relatório (restrita) & \textbullet~Preservar números do núcleo determinístico\newline \textbullet~Usar template se o texto LLM divergir & DEL, APR & \llmflag \\

\end{dtable}

\begin{callout}
\callouttitle{Serviços de plataforma não são agentes}
Comunicação A2A, gateway MCP, auditoria, registro de modelos, memória, SafetyPolicy e reserva de período são
\textbf{serviços compartilhados}: não têm objetivo local, iniciativa nem identidade social. Assim se evita a
\emph{agentificação excessiva}, audita-se o sistema como um todo e trocam-se peças (modelos, agentes, fontes)
sem efeito nas demais camadas. A \textbf{IA generativa} é \textbf{flag} (opcional, desligada por padrão) apenas em
Relatórios, Resultados e Orçamento, e só para narrativa. Números vêm sempre do núcleo determinístico e o registro
de modelos rejeita LLM em Orquestrador, Custos e Controles Internos (ADR-002).
\end{callout}

% =====================================================================
\section{Objetivos, percepções, ações e estado/conhecimento}
Para cada agente $i$: percepção $see_i: S\to P_i$, atualização de crenças $upd_i: B_i\times P_i\to B_i$ e decisão
$act_i: B_i\times K_i\to A_i$. Mantém-se a separação \textbf{Estado $\neq$ Conhecimento $\neq$ Memória}: o estado
muda a cada ciclo; o conhecimento (regras de rateio, alçadas, fórmulas) é versionado; a memória contém apenas saídas
\textbf{aprovadas} de ciclos anteriores (\code{SVC-MEM}). Em Controles Internos o ciclo é BDI explícito: crenças
sobre lançamentos $\rightarrow$ desejo de conformidade $\rightarrow$ intenção de aprovar, apontar ou vetar.

\begin{dtable}{L{2.35cm} Y Y Y Y}
\headrow \hd{Agente $\rightarrow$ objetivo} & \hd{Percepções (P)} & \hd{Estado interno (B)} & \hd{Conhecimento (K)} & \hd{Ações (A)}\\ \midrule\endhead
% ARQUIVO GERADO por tools/trace.py a partir de models/*.yaml — não editar à mão
\textbf{Orquestrador}\newline{\scriptsize\color{muted}AG-ORQ $\rightarrow$ G-01} & \textbullet~Solicitação do usuário (intenção, competência, entidade)\newline \textbullet~Mensagens agree/refuse/inform/failure dos especialistas\newline \textbullet~Parecer e veto de Controles Internos\newline \textbullet~Decisão do Approval Gate & \textbullet~DAG de subtarefas e status por conversation\_id\newline \textbullet~Reserva de período ativa\newline \textbullet~Relógio lógico e prazos reply\_by & \textbullet~Catálogo de intenções e planos (fechar mês, avaliar orçamento)\newline \textbullet~Diretório de papéis e Agent Cards & \textbullet~Enviar request a especialistas\newline \textbullet~Reservar e liberar período\newline \textbullet~Solicitar verificação de conformidade\newline \textbullet~Propor publicação ao Approval Gate\newline \textbullet~Entregar resposta final ou falha explícita \\
\textbf{Agente de Custos e Rentabilidade}\newline{\scriptsize\color{muted}AG-CUS $\rightarrow$ G-02} & \textbullet~Request de apuração (competência, entidade, dimensão)\newline \textbullet~Query-ref de Resultados\newline \textbullet~Balancete e lançamentos via MCP & \textbullet~Custos apurados do período e versão do rateio aplicado\newline \textbullet~Pendências de dados de origem & \textbullet~Critérios de rateio e classificação fixo/variável\newline \textbullet~Memória aprovada de ciclos anteriores & \textbullet~Calcular margem de contribuição e rentabilidade\newline \textbullet~Responder inform com proveniência\newline \textbullet~Responder failure com dado ausente \\
\textbf{Agente de Orçamento e Controle Orçamentário}\newline{\scriptsize\color{muted}AG-ORC $\rightarrow$ G-03} & \textbullet~Request de avaliação orçamentária\newline \textbullet~Query-ref de Resultados\newline \textbullet~Orçado e realizado via MCP & \textbullet~Desvios por conta e centro de custo\newline \textbullet~Forecast do trimestre e premissas & \textbullet~Limiar de desvio (\% e R\$) e método de forecast\newline \textbullet~Memória aprovada de justificativas de desvio & \textbullet~Sinalizar desvio acima do limiar\newline \textbullet~Projetar forecast\newline \textbullet~Propor alteração de orçamento (restrita) \\
\textbf{Agente de Resultados e Indicadores}\newline{\scriptsize\color{muted}AG-RES $\rightarrow$ G-04} & \textbullet~Request de consolidação de resultado\newline \textbullet~Inform de custos (CUS) e de orçamento (ORC)\newline \textbullet~Índices de mercado via MCP & \textbullet~DRE e KPIs em construção e pendências de insumos & \textbullet~Estrutura da DRE gerencial e fórmulas de KPIs\newline \textbullet~Memória aprovada de fechamentos anteriores & \textbullet~Enviar query-ref a CUS e ORC\newline \textbullet~Consolidar DRE, EBITDA, DFC e KPIs\newline \textbullet~Informar resultado ao Orquestrador \\
\textbf{Agente de Controles Internos}\newline{\scriptsize\color{muted}AG-CIN $\rightarrow$ G-05} & \textbullet~Request de verificação do pacote de fechamento\newline \textbullet~Lançamentos, aprovadores e conciliações via MCP\newline \textbullet~Waiver ou correção informados pelo Orquestrador & \textbullet~Crenças sobre conformidade de cada lançamento\newline \textbullet~Intenções ativas (aprovar, apontar, vetar) & \textbullet~Regras de alçada, SoD, duplicidade e competência (base normativa) & \textbullet~Emitir parecer conforme / não conforme\newline \textbullet~Vetar publicação por não conformidade crítica\newline \textbullet~Reavaliar após correção \\
\textbf{Agente de Relatórios Gerenciais e Análises}\newline{\scriptsize\color{muted}AG-REL $\rightarrow$ G-06} & \textbullet~Request de relatório com resultados consolidados e parecer de conformidade\newline \textbullet~Memória aprovada relevante & \textbullet~Rascunho do relatório e versão do template & \textbullet~Templates de relatório executivo\newline \textbullet~Memória aprovada (embeddings e grafo) & \textbullet~Gerar relatório por template\newline \textbullet~Gerar narrativa LLM (se flag ativa) e validar números\newline \textbullet~Propor publicação (restrita) \\

\end{dtable}

% =====================================================================
\section{Interações, coordenação e guardrails de segurança}

\subsection{Comunicação: A2A entre agentes, MCP com plataformas}
Toda comunicação entre agentes usa a mensagem imutável $m=\langle s, r, p, c, o, k\rangle$ (remetente, destinatário,
performativo FIPA-ACL \parencite{fipa2002acl}, conteúdo, ontologia \code{ctrl-fin-v1} e \code{conversation\_id}), com prazo lógico
\code{reply\_by}: respostas tardias são descartadas deterministicamente. \textbf{A2A é social} (agente $\leftrightarrow$
agente, com semântica de protocolo; \cite{a2a2025}); \textbf{MCP é instrumental} (agente $\rightarrow$ fonte de dados, como
\emph{tool}; \cite{mcp2025}). Nenhum agente usa MCP para conversar com outro.

\subsection{Mecanismos de coordenação}
\begin{dtable}{L{3.0cm} Y L{3.6cm}}
\headrow \hd{Dependência} & \hd{Situação no domínio} & \hd{Protocolo}\\ \midrule\endhead
Tarefa & Orquestrador delega subtarefas do DAG (padrão supervisor/especialistas) & \code{PROT-DEL} (FIPA-Request)\\ \rowsep
Informação & Resultados obtém custos e orçamento consolidados sem acessar regras alheias & \code{PROT-QRY} (Query-Ref)\\ \rowsep
Autoridade / norma & Publicação depende do parecer de Controles Internos, que pode \textbf{vetar} & \code{PROT-CMP} (verificação + veto)\\ \rowsep
Recurso exclusivo & Um fechamento por (entidade, competência), com TTL & \code{PROT-LCK} (reserva)\\ \rowsep
Autorização humana & Lançar ajuste no ERP, alterar orçamento, publicar relatório & \code{PROT-APR} (Approval Gate)\\ \rowsep
Fronteira & Pedido do usuário e resposta final explícita & \code{PROT-SOL} (FIPA-Request)\\
\end{dtable}

O Contract Net não é o mecanismo principal porque cada subtarefa tem um único papel competente: não há disputa a
leiloar. Ele fica registrado como evolução para o caso de várias instâncias do mesmo papel (por exemplo, um agente
de Custos por unidade de negócio).


\subsection{Protocolos como máquinas de estados}
Os \nProt{} protocolos são especificados como FSMs (\code{models/protocols.yaml}), seguindo a estrutura do FIPA-Request \parencite{fipa2002request}. As transições válidas originam
testes nominais e as \textbf{transições inválidas} originam testes negativos, a executar na DSM2
(Figura~\ref{fig:fsm}): uma sequência que o modelo proíbe deve ser rejeitada mesmo que cada mensagem isolada
seja bem formada.

\begin{figure}[!htbp]
  \centering
  \resizebox{\linewidth}{!}{% Figura editável (TikZ): FSMs de PROT-DEL e PROT-APR, com transições inválidas (origem de testes negativos).
% Fonte das transições: models/protocols.yaml (diagramas Mermaid gerados em diagrams/dsm1/fsm_*.mmd).
\begin{tikzpicture}[font=\sffamily\tiny, >={Stealth[length=1.6mm]},
  st/.style={draw=navy, fill=white, rounded corners=5pt, minimum height=0.5cm, minimum width=1.6cm, font=\sffamily\bfseries\tiny, text=navy},
  fin/.style={st, double, double distance=0.9pt},
  bad/.style={draw=bad, fill=bad!8, text=bad, rounded corners=5pt, minimum height=0.45cm, font=\sffamily\bfseries\tiny},
  tr/.style={->, navy, line width=0.5pt},
  inv/.style={->, bad, dashed, line width=0.6pt},
  l/.style={fill=white, inner sep=1pt, font=\sffamily\tiny},
]
% ---------------- PROT-DEL
\node[font=\sffamily\bfseries\scriptsize, text=accent, anchor=west] at (-0.8,2.05) {PROT-DEL · delegação de subtarefa};
\node[st]  (c)  at (0,0)       {CRIADA};
\node[st]  (ag) at (2.5,0)     {AGUARDANDO};
\node[st]  (ex) at (5.2,0)     {EM\_EXECUCAO};
\node[fin] (ok) at (7.7,0.65)  {CONCLUIDA};
\node[fin] (fa) at (7.7,-0.65) {FALHOU};
\node[fin] (re) at (2.5,-1.3)  {RECUSADA};
\node[fin] (xp) at (5.2,-1.3)  {EXPIRADA};
\node[bad] (v1) at (2.5,1.3)   {ProtocolViolation};
\draw[tr] (c) -- node[l, above]{request} (ag);
\draw[tr] (ag) -- node[l, above]{agree} (ex);
\draw[tr] (ex) -- node[l, above, sloped]{inform} (ok);
\draw[tr] (ex) -- node[l, below, sloped]{failure} (fa);
\draw[tr] (ag) -- node[l, left]{refuse} (re);
\draw[tr] (ag) -- node[l, below, sloped]{timeout} (xp);
\draw[tr] (ex) -- node[l, right]{timeout} (xp);
\draw[inv] (ag) -- node[l, left, text=bad]{inform antes de agree} (v1);

% ---------------- PROT-APR
\begin{scope}[xshift=10.0cm]
\node[font=\sffamily\bfseries\scriptsize, text=accent, anchor=west] at (-0.8,2.05) {PROT-APR · Approval Gate (HITL)};
\node[st]  (p)   at (0,0)       {PROPOSTA};
\node[st]  (pe)  at (2.6,0)     {PENDENTE};
\node[fin] (exe) at (5.5,0.65)  {EXECUTADA};
\node[fin] (rj)  at (5.5,-0.65) {REJEITADA};
\node[fin] (ng)  at (0,-1.3)    {NEGADA};
\node[fin] (ep)  at (2.6,-1.3)  {EXPIRADA};
\node[bad] (v2)  at (2.6,1.15)  {ApprovalRequired};
\draw[tr] (p) -- node[l, above]{require\_approval} (pe);
\draw[tr] (p.north) -- (0,1.65) -- node[l, above]{allow} (5.5,1.65) -- (exe.north);
\draw[tr] (p) -- node[l, left]{deny} (ng);
\draw[tr] (pe) -- node[l, above, sloped]{approve} (exe);
\draw[tr] (pe) -- node[l, below, sloped]{reject} (rj);
\draw[tr] (pe) -- node[l, right]{timeout} (ep);
\draw[inv] (pe) -- node[l, left, text=bad]{execute sem approve} (v2);
\end{scope}
\end{tikzpicture}
}
  \caption{FSMs de delegação e do Approval Gate. Estados com borda dupla são finais; em vermelho, transições inválidas que viram testes negativos.}
  \label{fig:fsm}
\end{figure}

\subsection{Cenário âncora: “feche o resultado do mês”}
A Figura~\ref{fig:seq} encadeia os protocolos no cenário âncora: reserva do período, subtarefas paralelas de Custos e
Orçamento, consultas A2A de Resultados, verificação com possibilidade de veto, relatório e publicação retida no
Approval Gate até a decisão humana. Cada mensagem carrega o mesmo \code{conversation\_id}, o que permite
reconstruir a cadeia causal a partir do trace.

\begin{figure}[!htbp]
  \centering
  \resizebox{\linewidth}{!}{% Figura editável (TikZ): cenário âncora "feche o resultado do mês" (fluxo nominal).
% Fonte equivalente em Mermaid: diagrams/dsm1/sequencia_fechar_mes.mmd
\begin{tikzpicture}[font=\sffamily\tiny, >={Stealth[length=1.6mm]},
  part/.style={draw=navy, fill=navy, text=white, rounded corners=2pt, minimum width=1.38cm, minimum height=0.5cm, font=\sffamily\bfseries\tiny, align=center},
  hpart/.style={part, fill=navy!8, text=navy},
  spart/.style={part, draw=plat, fill=platlight, text=plat},
  m/.style={->, navy, line width=0.5pt},
  r/.style={->, navy, dashed, line width=0.5pt},
  ml/.style={above, inner sep=1pt, fill=white, fill opacity=0.85, text opacity=1},
]
\def\H{8.3}
\foreach \x/\n/\s in {0/Solicitante/hpart, 1.75/AG-ORQ/part, 3.5/SVC-LCK/spart, 5.25/AG-CUS/part, 7.0/AG-ORC/part,
                      8.75/AG-RES/part, 10.5/AG-CIN/part, 12.25/AG-REL/part, 14.0/SVC-POL/spart, 15.75/Aprovador/hpart} {
  \node[\s] at (\x,0.35) {\n};
  \draw[muted!60, densely dashed] (\x,0.1) -- (\x,-\H);
}
% x: SOL 0 | ORQ 1.75 | LCK 3.5 | CUS 5.25 | ORC 7.0 | RES 8.75 | CIN 10.5 | REL 12.25 | POL 14.0 | APR 15.75
\newcommand{\M}[5][m]{\draw[#1] (#2,-#5*0.36) -- node[ml]{#4} (#3,-#5*0.36);}
\M{0}{1.75}{request(fechar\_mes)}{1}
\M{1.75}{3.5}{reserve}{2}
\M[r]{3.5}{1.75}{accept(ttl)}{3}
\M[r]{1.75}{0}{agree}{4}
\M{1.75}{5.25}{request(custos)}{5}
\M{1.75}{7.0}{request(orçado$\times$real)}{6}
\M[r]{5.25}{1.75}{inform(margens)}{7}
\M[r]{7.0}{1.75}{inform(desvios)}{8}
\M{1.75}{8.75}{request(consolidar)}{9}
\M{8.75}{5.25}{query-ref}{10}
\M[r]{5.25}{8.75}{inform}{11}
\M{8.75}{7.0}{query-ref}{12}
\M[r]{7.0}{8.75}{inform}{13}
\M[r]{8.75}{1.75}{inform(DRE, EBITDA, KPIs)}{14}
\M{1.75}{10.5}{request(conformidade)}{15}
\M[r]{10.5}{1.75}{inform(conforme)}{16}
\M{1.75}{12.25}{request(relatório)}{17}
\M[r]{12.25}{1.75}{inform(template | LLM validado)}{18}
\M{1.75}{14.0}{propose(publicar)}{19}
\M{14.0}{15.75}{REQUIRE\_APPROVAL}{20}
\M[r]{15.75}{14.0}{approve}{21}
\M{1.75}{3.5}{release}{22}
\M[r]{1.75}{0}{inform(resultado)}{22.6}
% bloco paralelo e alternativa de veto
\draw[accent, rounded corners=2pt] (1.55,-1.67) rectangle (7.35,-3.06);
\node[anchor=north east, font=\sffamily\bfseries\tiny, text=accent] at (1.5,-1.67) {par};
\draw[warn, rounded corners=2pt] (1.55,-5.27) rectangle (10.75,-5.94);
\node[anchor=north west, font=\sffamily\bfseries\tiny, text=warn, fill=white, inner sep=1pt] at (10.8,-5.31) {alt: veto $\rightarrow$ release + failure};
\node[anchor=east, font=\sffamily\tiny, text=muted] at (16.3,-\H-0.25) {toda mensagem é interceptada por SVC-AUD};
\end{tikzpicture}
}
  \caption{Cenário âncora “feche o resultado do mês” (fluxo nominal). Na alternativa de veto, o Orquestrador libera a reserva e responde \texttt{failure}.}
  \label{fig:seq}
\end{figure}

\subsection{Fronteira de autoridade e normas (human-in-the-loop)}
\begin{callout}
\callouttitle{Política de segurança: capacidade $\neq$ autorização}
Nenhuma proposta de agente, e muito menos de um LLM, dispara efeito externo sem passar pela
\textbf{SafetyPolicy} determinística, que classifica cada ação em três zonas:
\decisao{ALLOW} (leitura e cálculo internos);
\decisao{REQUIRE_APPROVAL} (ajuste no ERP, alteração de orçamento, publicação, retidos no Approval Gate até
decisão humana); \decisao{DENY} (fora da allowlist do papel; bloqueio fail-closed, sem opção de aprovação).
\end{callout}

\begin{dtable}{L{1.2cm} L{4.1cm} Y L{2.4cm} L{1.9cm}}
\headrow \hd{Norma} & \hd{Regra} & \hd{Consequência} & \hd{Decisão} & \hd{Imposta por}\\ \midrule\endhead
% ARQUIVO GERADO por tools/trace.py a partir de models/*.yaml — não editar à mão
\textbf{NORM-01} & Efeito externo exige aprovação humana & Ação retida no Approval Gate até approve; timeout expira sem executar & \decisao{REQUIRE_APPROVAL} & SVC-POL \\
\textbf{NORM-02} & Fora da allowlist do papel é negado & Bloqueio fail-closed, sem opção de aprovação, registrado em log & \decisao{DENY} & SVC-POL, SVC-MCP \\
\textbf{NORM-03} & Período em fechamento é exclusivo & reserve-failure explicado; reserva expira por TTL & \decisao{DENY} & SVC-LCK \\
\textbf{NORM-04} & Não conformidade crítica bloqueia publicação & Publicação vetada até correção reavaliada ou waiver humano justificado & \decisao{DENY} & SVC-POL \\
\textbf{NORM-05} & LLM não calcula números nem autoriza ações & Texto descartado e substituído pelo template; proposta tratada como não autorizada & \decisao{DENY} & SVC-MOD, SVC-POL \\
\textbf{NORM-06} & Logs e memória sem segredos nem dados pessoais desnecessários & Redaction; se a sanitização falhar, a gravação é recusada (fail-closed) & \decisao{OBRIGACAO} & SVC-AUD, SVC-MEM \\
\textbf{NORM-07} & LLM só por flag, só em papéis permitidos, dentro da cota & Chamada LLM bloqueada e agente segue em modo determinístico (template) & \decisao{DENY} & SVC-MOD \\
\textbf{NORM-08} & Somente saída aprovada vira memória & Não é indexada em embeddings nem no grafo & \decisao{DENY} & SVC-MEM \\

\end{dtable}

% =====================================================================
\section{Métricas de sucesso e matriz de rastreabilidade}\label{sec:trace}

\subsection{Métricas e evidências observáveis}
As métricas serão coletadas dos traces JSONL (\code{SVC-AUD}) em cenários sintéticos com gabarito e seeds
controlados (a partir da DSM2), sempre contra a baseline em workflow.

\begin{dtable}{L{1.0cm} L{3.3cm} Y L{3.0cm} L{2.1cm}}
\headrow \hd{ID} & \hd{Métrica} & \hd{Fórmula} & \hd{Alvo} & \hd{Mede}\\ \midrule\endhead
% ARQUIVO GERADO por tools/trace.py a partir de models/*.yaml — não editar à mão
\textbf{MET-01} & Taxa de conclusão de solicitações & solicitações com resposta final explícita (inform ou failure justificado) / total & >= 95\% nos cenários controlados & \idref{REQ-F01} \\ \rowsep
\textbf{MET-02} & Acurácia numérica contra gabarito & itens de DRE, margens e desvios com |valor - gabarito| <= R\$ 0,01 / total de itens & 100\% & \idref{REQ-F02}, \idref{REQ-F03}, \idref{REQ-F04} \\ \rowsep
\textbf{MET-03} & Detecção de inconformidades (precisão e recall) & recall e precisão sobre inconformidades plantadas por severidade & recall 100\% em críticas; precisão >= 90\% & \idref{REQ-F05} \\ \rowsep
\textbf{MET-04} & Overhead de comunicação & mensagens A2A por solicitação (comparado ao baseline workflow) & reportar e justificar; sem alvo fixo & \idref{REQ-F01}, \idref{REQ-F04} \\ \rowsep
\textbf{MET-05} & Contenção por guardrail & ações restritas bloqueadas antes do efeito sem aprovação / ações restritas tentadas & 100\% & \idref{REQ-S01}, \idref{REQ-S02} \\ \rowsep
\textbf{MET-06} & Consistência numérica do texto gerado & relatórios LLM com 100\% dos números iguais ao núcleo ou com fallback / relatórios LLM & 100\% & \idref{REQ-S03}, \idref{REQ-F06} \\ \rowsep
\textbf{MET-07} & Completude de trace & solicitações cuja cadeia causal é reconstruível pelo conversation\_id / total & 100\% & \idref{REQ-O01}, \idref{REQ-O02}, \idref{REQ-O04} \\ \rowsep
\textbf{MET-08} & Custo por solicitação & tokens de entrada/saída e custo estimado por solicitação, flag ligada vs desligada & custo zero com flag desligada; dentro da cota com flag ligada & \idref{REQ-O03}, \idref{REQ-C01} \\ \rowsep
\textbf{MET-09} & Latência do fechamento & tempo (lógico e de parede) da solicitação até a resposta final vs baseline & reportar distribuição em 10 seeds & \idref{REQ-F01}, \idref{REQ-NF01} \\ \rowsep
\textbf{MET-10} & Reaproveitamento de memória & análises que citam memória aprovada relevante / análises com memória disponível & >= 80\% no cenário de dois ciclos & \idref{REQ-F07} \\ \rowsep
\textbf{MET-11} & Vazamento em logs & ocorrências de campos proibidos (segredos, dados pessoais) nos arquivos de log e memória & 0 & \idref{REQ-P01} \\ \rowsep
\textbf{MET-12} & Reprodutibilidade & execuções com mesmo seed e mesmas versões que produzem o mesmo hash / execuções & 100\% & \idref{REQ-NF01}, \idref{REQ-NF02} \\

\end{dtable}

\subsection{Matriz de rastreabilidade ponta a ponta (AOSE)}
A cadeia liga cada necessidade do negócio a requisito, objetivo, papel/agente, protocolo/norma, componente, teste e
métrica, e a evidência do teste volta ao requisito (Figura~\ref{fig:cadeia}). Os \nObj{} objetos vêm de uma
\textbf{fonte única} (\code{models/*.yaml}). O script \code{tools/trace.py} falha diante de ID órfão, link quebrado
ou requisito sem objetivo, papel, protocolo/norma, teste ou métrica, e gera tanto esta tabela quanto a
documentação navegável: \textbf{\nReq/\nReq{} requisitos (100\%) com cadeia completa}. Os \nTest{} testes estão
\textbf{especificados, não executados} (nenhum é declarado aprovado antes de existir código). Componentes e
execuções são \nPend{} \emph{stubs} rastreáveis pendentes da DSM2.

\begin{figure}[H]
  \centering
  \resizebox{\linewidth}{!}{% Figura editável (TikZ): cadeia de rastreabilidade com exemplo concreto (REQ-S01).
\begin{tikzpicture}[font=\sffamily\scriptsize, >={Stealth[length=2mm]},
  box/.style={draw=navy, fill=white, rounded corners=3pt, minimum width=1.95cm, minimum height=0.95cm, align=center, line width=0.6pt},
  pend/.style={box, dashed, draw=muted, text=muted},
  ex/.style={font=\sffamily\bfseries\tiny, text=accent, align=center, anchor=north},
]
\node[box]  (n0) at (0,0)     {\textbf{Problema}};
\node[box]  (n1) at (2.38,0)  {\textbf{Requisito}};
\node[box]  (n2) at (4.76,0)  {\textbf{Objetivo}};
\node[box]  (n3) at (7.14,0)  {\textbf{Papel /}\\\textbf{Agente}};
\node[box]  (n4) at (9.52,0)  {\textbf{Protocolo /}\\\textbf{Norma}};
\node[pend] (n5) at (11.9,0)  {\textbf{Componente}\\(DSM2)};
\node[pend] (n6) at (14.28,0) {\textbf{Teste}\\(planejado)};
\node[ex] at (n0.south) {PROB-08};
\node[ex] at (n1.south) {REQ-S01};
\node[ex] at (n2.south) {G-09};
\node[ex] at (n3.south) {ROLE-APR · AG-ORQ};
\node[ex] at (n4.south) {PROT-APR · NORM-01};
\node[ex] at (n5.south) {COMP-POL};
\node[ex] at (n6.south) {TEST-12};
\draw[->, navy, thick] (n0) -- (n1);
\draw[->, navy, thick] (n1) -- (n2);
\draw[->, navy, thick] (n2) -- (n3);
\draw[->, navy, thick] (n3) -- (n4);
\draw[->, navy, thick] (n4) -- (n5);
\draw[->, navy, thick] (n5) -- (n6);
\draw[->, accent, dashed, thick] (n6.north) -- ++(0,0.45) -| node[pos=0.25, above, font=\sffamily\tiny, text=accent] {evidência do teste retroalimenta o critério de aceite} (n1.north);
\node[box, fill=accentlight, minimum width=1.6cm] (met) at (2.38,-1.8) {\textbf{Métrica}};
\node[font=\sffamily\bfseries\tiny, text=accent, anchor=west] at (met.east) {MET-05: 100\% das ações restritas contidas antes do efeito};
\draw[->, accent, thick] (met) -- (2.38,-0.95);
\end{tikzpicture}
}
  \caption{Cadeia de rastreabilidade com exemplo concreto. Em tracejado, itens que se materializam na DSM2.}
  \label{fig:cadeia}
\end{figure}

\begin{dtable}{L{1.35cm} L{1.35cm} L{1.0cm} L{2.6cm} Y L{3.3cm} L{2.2cm}}
\headrow \hd{Problema} & \hd{Requisito} & \hd{Objetivo} & \hd{Agente} & \hd{Protocolo / Norma} & \hd{Teste (planejado)} & \hd{Métrica}\\ \midrule\endhead
% ARQUIVO GERADO por tools/trace.py a partir de models/*.yaml — não editar à mão
PROB-01 & \idref{REQ-F01} & G-01 & ORQ & SOL, DEL & TEST-01, TEST-02, TEST-03 & MET-01, MET-04, MET-09 \\
PROB-02 & \idref{REQ-F02} & G-02 & CUS & DEL & TEST-04 & MET-02 \\
PROB-03 & \idref{REQ-F03} & G-03 & ORC & DEL, QRY & TEST-05 & MET-02 \\
PROB-01 & \idref{REQ-F04} & G-04 & CUS, ORC, RES & QRY, DEL & TEST-06, TEST-07 & MET-02, MET-04 \\
PROB-04 & \idref{REQ-F05} & G-05 & CIN, ORQ & CMP, NORM-04 & TEST-08, TEST-09 & MET-03 \\
PROB-05 & \idref{REQ-F06} & G-06 & REL & DEL, NORM-05, NORM-07 & TEST-10, TEST-15 & MET-06 \\
PROB-07 & \idref{REQ-F07} & G-08 & CUS, ORC, REL, RES & APR, NORM-08 & TEST-11 & MET-10 \\
PROB-08 & \idref{REQ-S01} & G-09 & ORC, ORQ, REL & APR, NORM-01, NORM-02 & TEST-12, TEST-13 & MET-05 \\
PROB-01, PROB-04 & \idref{REQ-S02} & G-01, G-05 & ORQ & LCK, NORM-03 & TEST-14 & MET-05 \\
PROB-05, PROB-06 & \idref{REQ-S03} & G-06, G-09 & ORC, REL, RES & APR, NORM-05 & TEST-15 & MET-06 \\
PROB-06 & \idref{REQ-O01} & G-07 & CIN, CUS, ORC, ORQ, REL, RES & DEL, NORM-06 & TEST-16 & MET-07 \\
PROB-06 & \idref{REQ-O02} & G-07 & ORC, REL, RES & NORM-06, NORM-07 & TEST-17 & MET-07 \\
PROB-06 & \idref{REQ-O03} & G-07 & ORC, REL, RES & NORM-07 & TEST-18, TEST-22 & MET-08 \\
PROB-06 & \idref{REQ-O04} & G-07 & APR, SOL & NORM-06 & TEST-19 & MET-07 \\
PROB-06 & \idref{REQ-C01} & G-07, G-09 & ORC, REL, RES & NORM-07 & TEST-20, TEST-21, TEST-22 & MET-08 \\
PROB-06 & \idref{REQ-P01} & G-07 & CIN, ORQ, REL & NORM-06 & TEST-23 & MET-11 \\
PROB-06 & \idref{REQ-NF01} & G-07 & ORQ, REL & DEL, NORM-07 & TEST-10, TEST-24 & MET-09, MET-12 \\
PROB-06 & \idref{REQ-NF02} & G-07 & CIN, REL & NORM-07 & TEST-25 & MET-12 \\

\end{dtable}

\begin{callout}
\callouttitle{Disposições finais e próximos passos}
A DSM1 entrega cenário, fronteira, requisitos, organização de papéis, modelo PEAS/BDI, protocolos, normas e
métricas com rastreabilidade validada automaticamente. Os modelos editáveis estão em \code{models/*.yaml},
\code{diagrams/dsm1/*.mmd} (Mermaid) e \code{entregas/dsm1/figures/*.tex} (TikZ). Na organização dos arquivos e da
documentação foi utilizado o modelo de IA \textbf{Claude Opus 5.5} (Anthropic), conforme \code{USO\_DE\_IA.md}. A \textbf{DSM2} (arquitetura e protótipo) está apenas pré-estruturada: runtime, transporte
A2A concreto, componentes e execução dos testes serão decididos e documentados nela.\par\vspace{2pt}
{\RaggedRight\textbf{Documentação navegável:} \href{\docsURL}{\texttt{\docsURL}}\\
\textbf{Repositório:} \href{\repoURL}{\texttt{\repoURL}}\par}
\end{callout}

\clearpage
\begingroup\small\RaggedRight
\printbibliography[heading=bibintoc, title={Referências}]
\endgroup

\end{document}
